\documentclass[10pt]{article}

\usepackage[margin=0.9in]{geometry}
\usepackage{lmodern}
\usepackage[T1]{fontenc}
\usepackage{microtype}
\usepackage{amsmath,amssymb}
\usepackage{stmaryrd}
\usepackage{pifont}
\usepackage{mathpartir}
\usepackage{booktabs}
\usepackage{enumitem}
\usepackage{xcolor}
\usepackage[colorlinks=true,linkcolor=blue!50!black,urlcolor=blue!50!black]{hyperref}

\setlist{nosep,leftmargin=1.4em}

% ---------------------------------------------------------------- notation
\newcommand{\ov}[2]{\{\nu{:}#1 \mid #2\}}          % demonic / over
\newcommand{\un}[2]{[\nu{:}#1 \mid #2]}            % angelic / under
\newcommand{\Prec}[1]{\mathsf{P}(#1)}               % precise: {nu = e} /\ [nu = e]
\newcommand{\Eq}[1]{\mathsf{Eq}(#1)}                % {nu = v} /\ [nu = v]
\newcommand{\Nat}{\mathsf{nat}}
\newcommand{\Bool}{\mathsf{bool}}
\newcommand{\List}{\mathsf{list}}
\newcommand{\Tree}{\mathsf{tree}}
\newcommand{\wand}{\mathrel{-\!\!*}}
\newcommand{\inter}{\sqcap}
\newcommand{\Sum}{\oplus}
\newcommand{\pers}{\Box}
\newcommand{\entc}{\mathrel{,}}
\newcommand{\sep}{\mathrel{*}}
\newcommand{\vd}{\vdash_{\!x}}
\newcommand{\den}[1]{\llbracket #1 \rrbracket}
\newcommand{\ctxsub}[1]{\mathrel{\preceq_{#1}}}
\newcommand{\ctxcase}[1]{\mathrel{\preceq^{\vee}_{#1}}}
\newcommand{\sub}{\mathrel{<:}}
\newcommand{\kw}[1]{\ensuremath{\mathbf{#1}}}
\newcommand{\rocq}[1]{\texttt{#1}}
\newcommand{\RN}[1]{\textsc{#1}}
\newcommand{\obl}[1]{\ensuremath{\textcolor{red!60!black}{\scriptstyle\langle #1\rangle}}}
\newcommand{\qed}[1]{\ensuremath{\textcolor{green!40!black}{\scriptstyle\langle #1\rangle}}}

\title{Typing the Case Studies of \emph{Context Types} in Rocq:\\
Extended Rules, Concrete Derivations, and Open Issues}
\author{Summary of the \texttt{ProofAutomation} / \texttt{CaseStudy} development}
\date{October 2026}

\begin{document}
\maketitle

\begin{abstract}
We report on mechanizing typing derivations for the case studies of the draft
paper \emph{Context Types: Bunched Modalities for Unifying Safety and
Reachability}. The Rocq development already proves the Fundamental theorem for
the core judgment; the case studies, however, only had basic (erased) typings.
We added derivations in a case-study judgment $\Sigma;\Gamma \vd e : \tau$,
extended it with four rule families (a dependent variable rule, a
capability case split, fixpoint unrolling, and paper-form constructor matches),
and proved each new rule sound against the core denotation. We then derived
the paper's examples, found several places where the paper's stated type is
semantically false or vacuous, and isolated one remaining design issue in
\RN{T-Fix}. All structural steps of all derivations are checked (\rocq{Qed});
semantic side conditions are isolated as named obligations.
\end{abstract}

% =====================================================================
\section{Setting and trust model}

\paragraph{Judgment.} \rocq{ProofAutomation/ExtendedTyping.v} defines
$\Sigma;\Gamma\vd e:\tau$ (\rocq{has\_extended\_context\_type}). It contains
every rule of the core judgment \rocq{has\_context\_type}
(\rocq{typing\_to\_extended\_typing}), plus case-study conveniences. Types are
$\ov{b}{\phi}$ (demonic/over), $\un{b}{\phi}$ (angelic/under), $\tau_1\inter\tau_2$,
$\tau_1\Sum\tau_2$, $x{:}\tau_x\to\tau$, $x{:}\tau_x\wand\tau$ and $\pers\tau$;
contexts are built with entangled ``$\entc$'', independent ``$\sep$'' and sum
``$\Sum$''. We write $\Prec{e}$ for the precise type
$\ov{b}{\nu = e}\inter\un{b}{\nu = e}$ that the operator rules produce, and
$\Eq{v}$ for the match-equation type $\ov{b}{\nu=v}\inter\un{b}{\nu=v}$.
Paper \textsf{int} is translated to Rocq $\Nat$ (truncated subtraction).

\paragraph{Semantic side conditions.} \RN{T-Sub} needs
$\Gamma\vdash\tau_1\sub\tau_2$ (\rocq{sub\_type\_under}: $\den{\Gamma}\models\den{\tau_1}e\Rightarrow\den{\tau_2}e$
for all well-typed $e$), and \RN{T-CtxSub} needs
$\Gamma_1\ctxsub{X}\Gamma_2$ (\rocq{ctx\_sub\_under}: every
$r\models\den{\Gamma_1}$ has some $r'\models\den{\Gamma_2}$ with $r|_X\sqsubseteq r'$).

\paragraph{Trust model.} Every derivation step is checked by Rocq (\rocq{Qed}).
Each semantic side condition is a separately named lemma
\rocq{<program>\_obl\_<what>}. Where an obligation is not machine-proved it is
\rocq{Admitted} together with an informal argument against the Rocq denotation.
Every case-study file was audited three times by adversarial reviewers who
tried to build counterexample capabilities; none of the remaining obligations
was refuted. The infrastructure (\rocq{ExtendedTyping.v},
\rocq{ExtendedSubsumption.v}, \rocq{ExtendedRuleSoundness.v},
\rocq{CaseTypes.v}) contains no \rocq{Admitted}. The project builds with
\rocq{make}, and no file declares an \rocq{Axiom}.

\begin{center}\small
\begin{tabular}{lrr@{\qquad}lrr}
\toprule
File & \rocq{Qed} & \rocq{Admitted} & File & \rocq{Qed} & \rocq{Admitted}\\
\midrule
\rocq{Section1Intro.v} & 56 & 16 & \rocq{MayMust.v} & 84 & 23\\
\rocq{Section2Typing.v} & 138 & 32 & \rocq{BSTGenerator.v} & 116 & 24\\
\rocq{Section3Logic.v} & 50 & 9 & \rocq{DisproveMerge.v} & 96 & 13\\
\rocq{Section4Extension.v} & 79 & 12 & \rocq{ExtendedRuleSoundness.v} & 189 & 0\\
\bottomrule
\end{tabular}
\end{center}

% =====================================================================
\section{Corrections to the case-study infrastructure}

\paragraph{Angelic top.} The aliases $U[b]$ used to be
\rocq{under\_ty b qual\_top}. Because \rocq{qual\_top} mentions no logic
variable, its exact atom ranges over the empty set and $\un{b}{\top}$ became
the \emph{trivial} under type, inhabited by every total term (proved:
$\ov{b}{q}\sub U[b]$ for any $q$). The paper's $\un{b}{\top}$ means ``every value of
$b$ is reachable''. The aliases now use
$\textsf{q\_base\_all}\;b \triangleq \nu \in b$, a qualifier that mentions $\nu$.
The demonic top $\ov{b}{\top}$ stays trivial, as intended. The same subtlety
exists on paper: under Def.~4.7, $\mathsf{Atom}(\top)$ has $\mathsf{FV}=\emptyset$.

\paragraph{Generic subtyping lemmas.} \rocq{ExtendedSubsumption.v} now proves
transitivity, the $\inter$-projections and commutation, over-closure by
qualifier entailment, and vacuous subtyping from an unsatisfiable context, all
with \rocq{Qed}. Three previously admitted lemmas were \emph{false} and were
replaced by corrected, proved versions:
\begin{itemize}
\item $\un{b}{q_1}\sub\un{b}{q_2}$ from $q_2\Rightarrow q_1$ fails if $q_2$
  is unsatisfiable, because $\mathsf{FUnder}(\mathsf{FAtom}\,\bot)$ has no
  model when worlds are non-empty. It also fails if $q_2$ mentions $\nu$
  but $q_1$ does not. The corrected lemma adds $\mathsf{vars}(q_2)\subseteq\mathsf{vars}(q_1)$
  and $\exists s.\,q_2(s)$.
\item $\ov{b}{q_o}\inter\un{b}{q_u}\sub\un{b}{q_o\wedge q_u}$ fails for
  $q_u=\top$. The corrected lemma adds $\mathsf{vars}(q_o)\subseteq\mathsf{vars}(q_u)$.
\end{itemize}
New lemmas proved along the way include $\tau_x\to\tau\ \sub\ \tau_x\wand\tau$ in
every context, and closure facts for $\Sum$: for example
$\ov{b}{q_1}\Sum\un{b}{q}\sub\un{b}{q}$.

% =====================================================================
\section{New rules}

Figure~\ref{fig:rules} lists the rules added to $\vd$. Each is justified by a
theorem in \rocq{ExtendedRuleSoundness.v} of the same shape as the core
Fundamental theorem: if the premises are semantically valid
($\den{\Gamma_i}\models\den{\tau_i}e_i$) and the side conditions hold, then
$\den{\Gamma}\models\den{\tau}e$.

\begin{figure}[t]
\begin{mathpar}
\inferrule*[right=T-VarComma]
  {\Sigma \vdash^{\mathsf{WF}} \Gamma\entc x{:}\tau \\ \mathsf{fv}(\tau)\subseteq\mathsf{dom}\lfloor\Gamma\rfloor}
  {\Sigma;\ \Gamma \entc x{:}\tau \vd x : \tau}

\inferrule*[right=T-CtxCase]
  {\Sigma;\Gamma_a \vd e:\tau \\ \Sigma;\Gamma_b \vd e:\tau \\
   \Gamma \ctxcase{\mathsf{fv}(e)\cup\mathsf{fv}(\tau)} \Gamma_a \mid \Gamma_b}
  {\Sigma;\Gamma \vd e:\tau}
\\
\inferrule*[right=T-FixUnroll]
  {\mathsf{lc}(e_b) \\
   \Sigma;\Gamma \vd \lambda a{:}s.\ e_b[f \mapsto \kw{fix}\,f.\lambda a.e_b] : \tau}
  {\Sigma;\Gamma \vd \kw{fix}\,f.\lambda a{:}s.\ e_b : \tau}

\inferrule*[right=T-ListMatchCons]
  {\Sigma;\Gamma \vd h :: t : \Eq{x} \\ \Sigma;\Gamma \vd e_c : \tau}
  {\Sigma;\Gamma \vd \kw{match}\ x\ \kw{with}\ [\,]\to e_n \mid h::t \to e_c : \tau}
\\
\inferrule*[right=T-ListMatchBoth]
  {\Sigma;\Gamma_n \vd [\,] : \Eq{x} \\ \Sigma;\Gamma_c \vd h::t : \Eq{x} \\
   \Sigma;\Gamma_n \vd e_n:\tau_n \\ \Sigma;\Gamma_c \vd e_c:\tau_c}
  {\Sigma;\Gamma_n\Sum\Gamma_c \vd \kw{match}\ x\ \kw{with}\ [\,]\to e_n \mid h::t\to e_c : \tau_n\Sum\tau_c}
\end{mathpar}
\vspace{-1ex}
\[
\Gamma \ctxcase{X} \Gamma_a\mid\Gamma_b \;\triangleq\;
\forall r\models\den{\Gamma}.\ \exists r'.\ r|_X\sqsubseteq r' \wedge
\big(r'\models\den{\Gamma_a} \vee r'\models\den{\Gamma_b}\big)
\]
\caption{Rules added to the case-study judgment. The side condition of
T-CtxCase also requires well-formedness and erasure agreement on $X$. The tree-match rules are
analogous to the list ones, with $[\,]\mapsto\mathsf{Leaf}$, the pattern
$h::t$ replaced by $\mathsf{Node}(r,l,rr)$, and the extra premise $l\neq rr$.
The Rocq rules use named binders (\rocq{lam\_named}, \rocq{fixfun\_named},
\rocq{match\_list\_named}).}
\label{fig:rules}
\end{figure}

\paragraph{T-VarComma} (\rocq{EXT\_VarComma}; sound: \rocq{ext\_var\_comma\_sound}).
Core \RN{T-Var} types $x$ only in the singleton context $x{:}\tau$. The paper
checks such a context under the basic context $\Sigma,\lfloor\Gamma\rfloor$
(\textsc{WfComma}), so $\tau$ may mention $\Gamma$. In Rocq, $\Sigma$ is a
fixed parameter of the judgment and type well-formedness is checked against
$\lfloor\Gamma\rfloor$ only, so T-CtxSub alone cannot project
$\Gamma\entc x{:}\tau$ down to $x{:}\tau$ when $\tau$ depends on $\Gamma$.
T-VarComma keeps the dependencies in the context; T-CtxSub does all other
structural work (weakening, exchange, framing). It is what makes recursive
calls such as $f\,(x-1)$, whose argument needs $\nu\prec x$, and variables with
dependent types such as $y{:}\ov{\Nat}{\nu=x+1}$ usable at all.

\paragraph{T-CtxCase} (\rocq{EXT\_CtxCase}; sound: \rocq{ext\_ctx\_case\_sound}).
This is a disjunctive form of T-CtxSub: every capability of $\Gamma$, projected
on the free variables, is a capability of $\Gamma_a$ or of $\Gamma_b$. It
resolves the main expressiveness gap we found. The match rules require either
a split $\Gamma_t\Sum\Gamma_f$ into two \emph{non-empty} summands (both
branches reachable) or a uniform guard (one branch only). But a demonic
binding $x{:}\ov{\Nat}{\top}$ admits capabilities that reach only one branch
and capabilities that reach both, so no single choice works. For instance,
$x{:}\ov{\Nat}{\top}\vdash \kw{if}\ x{=}0\ \kw{then}\ 3\ \kw{else}\ 5 : \ov{\Nat}{\nu{=}3\vee\nu{=}5}$
was not derivable before. With T-CtxCase the derivation splits the
capabilities into ``true only'', ``false only'' and ``both'', and each case
uses the existing match rules.

\paragraph{T-FixUnroll} (\rocq{EXT\_FixUnrollNamed}; sound:
\rocq{ext\_fix\_unroll\_named\_sound\_ext}, for \emph{every} target type). A
fixpoint may be typed by its one-step unfolding. This is the standard
unrolling of incorrectness reasoning, used for the merge disproof
(\S\ref{sec:merge}). The proof is a world-level value transport: two
locally closed values with equal free variables whose applications are
equivalent have the same type denotation (\rocq{ext\_vt\_main}). The premise
$\mathsf{lc}(e_b)$ is necessary: a dangling bound variable would be captured
by different binders in the fixpoint and in its unfolding.

\paragraph{Paper-form matches} (sound: \rocq{ext\_list\_match\_*\_paper\_rule\_sound},
\rocq{ext\_tree\_match\_*}). As in the paper's \RN{T-Match}, the context of a
reachable branch binds the pattern variables itself, with types chosen by the
user. They are introduced beforehand by T-CtxSub, whose extension condition is
the semantic content. The pre-existing cofinite rules forced the pattern
binders to $\ov{\Nat}{\top}$ and $\ov{\List}{\top}$, which makes the equation
premise $h::t:\Eq{x}$ unprovable.

\paragraph{An unsound first attempt.} Our first version checked the branch
premises under $\Gamma\entc\Gamma_{\mathit{pat}}$ but concluded at $\Gamma$.
That is unsound, because an unmodelable $\Gamma_{\mathit{pat}}$ makes every
premise vacuous. Rocq derives
$\vd \kw{match}\ [1]\ \ldots : \ov{\Nat}{\bot}$ from that rule shape
(\rocq{ext\_list\_match\_cons\_only\_named\_unsound}, kept as a
machine-checked warning).

\paragraph{Named-binder variants.} We also added \RN{T-LamD}, \RN{T-LetD} and
\RN{T-Fix} variants with concrete binders
(\rocq{EXT\_LamDNamedErased}, \rocq{EXT\_LetDNamed}, \rocq{EXT\_FixNamed}). They
mirror the cofinite core rules in the same way as the pre-existing
\rocq{EXT\_LamNamed} and \rocq{EXT\_LetNamed}; their soundness is not proved
separately (\S\ref{sec:open}).

% =====================================================================
\section{Concrete derivations}

\subsection{Figure 6(a): may-analysis helps must-analysis}\label{sec:fig6a}

The program is
$f = \lambda g.\lambda x.\lambda y.\ \kw{let}\ c = y > 10\ \kw{in}\
\kw{if}\ c\ \kw{then}\ (\kw{let}\ gx = g\,x\ \kw{in}\ gx + 1)\ \kw{else}\ x + 1$,
with may-summary $\tau_g \triangleq \ov{\Nat}{\top}\to\ov{\Nat}{\nu\neq 0}$.
We formalized the three settings proposed by the authors.

\paragraph{Setting (1)} $x{:}\un{\Nat}{\nu=0}\to y{:}\un{\Nat}{\nu\le 10}\to\un{\Nat}{\nu=1}$
is \emph{semantically false}. Because ``$\to$'' entangles $y$ with $x$, the
capability
\[ r = \{(x{=}0,\,y{=}11)\}\ \cup\ \{(x{=}1,\,y{=}n)\mid n\le 10\}\]
satisfies both bindings: $x=0$ is reachable, and every $y\le10$ is reachable,
since $y$'s qualifier does not mention $x$ and is therefore not fibered by $x$.
Yet $f$ never returns $1$: the store $(0,11)$ takes the then-branch, where
$g\,0+1\neq1$, and every other store returns $2$. Matching on $y$ is exactly
what ``loses'' $x$'s guarantee. No rule can derive this type, since the rules
are sound.

\paragraph{Setting (2)} $x{:}\un{\Nat}{\nu=0}\to y{:}\un{\Nat}{\nu\le 10}\wand\un{\Nat}{\nu=1}$
is valid and derived (\rocq{fig6a\_2\_typing}). After the three abstractions
and the \kw{let} of $c$, the body context is
\[\Gamma_2 = (g{:}\tau_g\entc x{:}\un{\Nat}{\nu{=}0})\sep y{:}\un{\Nat}{\nu{\le}10}\ \entc\ c{:}\Prec{y>10}.\]
The derivation splits on whether $y$'s capability reaches a value $>10$. The
branch contexts are
$\Gamma_T = g{:}\tau_g\entc x{:}\ov{\Nat}{\top}\entc c{:}\Prec{\mathsf{true}}$ and
$\Gamma_F = g{:}\tau_g\entc x{:}\un{\Nat}{\nu{=}0}\entc c{:}\Prec{\mathsf{false}}$.
\[
\inferrule*[right=CtxCase]
 {\mathcal{D}_1 \\ \mathcal{D}_2 \\ \obl{\Gamma_2 \ctxcase{} \Gamma_F \mid \Gamma_T\Sum\Gamma_F}}
 {\Gamma_2 \vd \kw{if}\ c\ \kw{then}\ \ldots\ \kw{else}\ \ldots : \un{\Nat}{\nu{=}1}}
\qquad
\mathcal{D}_1 =
\inferrule*[right=MatchFalseOnly]
  {\Gamma_F\vd c:\Prec{\mathsf{false}} \\ \mathcal{D}_{\mathit{else}}}
  {\Gamma_F \vd \kw{if}\ldots : \un{\Nat}{\nu{=}1}}
\]
\[
\mathcal{D}_2 =
\inferrule*[right=Sub]
 {\inferrule*[right=MatchBoth]
    {\Gamma_T\vd c:\Prec{\mathsf{true}}\\ \Gamma_F\vd c:\Prec{\mathsf{false}}\\
     \mathcal{D}_{\mathit{then}} \\ \mathcal{D}_{\mathit{else}}}
    {\Gamma_T\Sum\Gamma_F\vd \kw{if}\ldots : \ov{\Nat}{\nu{\ne}1}\Sum\un{\Nat}{\nu{=}1}}
  \\ \qed{\Sum\text{-absorption}}}
 {\Gamma_T\Sum\Gamma_F \vd \kw{if}\ldots : \un{\Nat}{\nu{=}1}}
\]
\[
\mathcal{D}_{\mathit{else}} =
\inferrule*[right=Sub]
 {\inferrule*[right=BinOp]{ }{\Gamma_F\vd x+1:\Prec{x+1}} \\ \obl{\mathit{else\_eq1\_x0}}}
 {\Gamma_F\vd x+1:\un{\Nat}{\nu{=}1}}
\]
\[
\mathcal{D}_{\mathit{then}} =
\inferrule*[right=Let]
 {\inferrule*[right=AppFun]{\ldots}{\Gamma_T\vd g\,x:\ov{\Nat}{\nu{\ne}0}} \\
  \inferrule*[right=Sub]{\Gamma_T'\vd gx+1:\Prec{gx+1} \\ \obl{\mathit{then\_ne1}}}
   {\Gamma_T'\vd gx+1:\ov{\Nat}{\nu{\ne}1}}}
 {\Gamma_T\vd \kw{let}\ gx = g\,x\ \kw{in}\ gx+1 : \ov{\Nat}{\nu{\ne}1}}
\]
where $\Gamma_T' = \Gamma_T\entc gx{:}\ov{\Nat}{\nu\ne0}$.
The case-split obligation is where ``$\sep$'' matters. In the second case,
every store of the $g,x$ part is paired with some $y>10$ (making $c$ true)
\emph{and} with $y=0$ (making $c$ false). So $x=0$ is still reachable in the
$c=\mathsf{false}$ summand, which is what $\mathcal{D}_{\mathit{else}}$ needs.
With ``$\entc$'' this pairing fails: the counterexample of setting (1) is
exactly such an unpaired capability. The may-summary of $g$ enters only
through \textit{then\_ne1}: the then-branch never produces $1$, so it can be
absorbed into the angelic $\un{\Nat}{\nu=1}$.

\paragraph{Setting (3)}
$x{:}\un{\Nat}{\top}\to y{:}\un{\Nat}{(x{=}0\Rightarrow\nu{=}12)\wedge(x{>}0\Rightarrow\nu{\le}10)}\to\tau_3$.
The authors left the result type open. The strongest sound one is
$\tau_3 = \ov{\Nat}{\nu\neq1}\Sum\un{\Nat}{\nu\geq2}$, derived as
\rocq{fig6a\_3\_typing}. Both branches are reachable: $x=0$ forces $y=12$,
taking the then-branch, which never returns $1$ by the may-summary of $g$. In
the else-branch, every $x\ge1$ is paired with $y=0$, so $x+1$ reaches every
$n\ge2$. A result type mentioning $\un{\Nat}{\nu=1}$ is unsound, because $x=0$
is paired only with $y=12$.

\paragraph{Figure 6(b).} With the paper's summary
$g:(\ov{\Nat}{\top}\to\ov{\Nat}{\nu>0})\inter(\un{\Nat}{\top}\to\un{\Nat}{\nu>0})$,
the stated type $x{:}\ov{\Nat}{\top}\to\ov{\Nat}{\nu{=}2}\Sum\ov{\Nat}{\nu{=}5}$
is false. The function $g=\lambda n.n+1$ satisfies both summaries, yet the
singleton capability $x=0$ yields only $2$, and ``$\Sum$'' needs both parts
non-empty. The angelic version $x{:}\un{\Nat}{\top}\to\ldots$ is derived. The
safety-only type $x{:}\ov{\Nat}{\top}\to\ov{\Nat}{\nu\geq0}$ is derived with
two nested T-CtxCase splits on $a>10$.

\subsection{A demonic guard: the introduction's Peano subtraction}\label{sec:minus}

For $\lambda x.\lambda y.\ \kw{let}\ c = x\le y\ \kw{in}\ \kw{if}\ c\ \kw{then}\ 0\ \kw{else}\ x-y$
at $\tau_a = x{:}\ov{\Nat}{\nu>0}\to y{:}\ov{\Nat}{\nu>0}\to\ov{\Nat}{\nu\le x}$
(\rocq{nat\_minus\_program\_tau\_a}), the guard is not determined by the
demonic context
$\Gamma_1 = x{:}\ov{\Nat}{\nu{>}0}\entc y{:}\ov{\Nat}{\nu{>}0}\entc c{:}\Prec{x\le y}$.
A context cannot directly say ``the guard is not constant'', so we nest two
splits. The intermediate context
$\Gamma_{\neg T} = \ldots\entc c{:}\un{\Bool}{\nu=\mathsf{false}}$ states
``$\mathsf{false}$ is reachable''. Its type is a closed under type, so no
fibering is involved. Let $\Gamma_T$ and $\Gamma_F$ rebind $c$ to
$\Prec{\mathsf{true}}$ and $\Prec{\mathsf{false}}$, and let
$\mathit{if}$ abbreviate the conditional.
Write $\tau_{\le} = \ov{\Nat}{\nu\le x}$. The three branch derivations are
\[
\mathcal{D}_T =
\inferrule*[right=MatchTrueOnly]
 {\Gamma_T\vd c:\Prec{\mathsf{true}} \\
  \inferrule*[right=Sub]{\Gamma_T\vd 0:\Prec{0} \\ \obl{zero}}{\Gamma_T\vd 0:\tau_{\le}}}
 {\Gamma_T\vd \mathit{if}:\tau_{\le}}
\]
\[
\mathcal{D}_F =
\inferrule*[right=MatchFalseOnly]
 {\Gamma_F\vd c:\Prec{\mathsf{false}} \\
  \inferrule*[right=Sub]{\Gamma_F\vd x-y:\Prec{x-y} \\ \obl{false\_result}}{\Gamma_F\vd x-y:\tau_{\le}}}
 {\Gamma_F\vd \mathit{if}:\tau_{\le}}
\]
\[
\mathcal{D}_{TF} =
\inferrule*[right=Sub]
 {\inferrule*[right=MatchBoth]
   {\Gamma_T\vd c:\Prec{\mathsf{true}} \\ \Gamma_F\vd c:\Prec{\mathsf{false}} \\ \ldots}
   {\Gamma_T\Sum\Gamma_F\vd \mathit{if}:\tau_{\le}\Sum\tau_{\le}}
  \\ \obl{sum}}
 {\Gamma_T\Sum\Gamma_F\vd \mathit{if}:\tau_{\le}}
\]
and they are combined by two case splits:
\[
\inferrule*[right=CtxCase]
 {\mathcal{D}_T \\
  \inferrule*[right=CtxCase]
   {\mathcal{D}_F \\ \mathcal{D}_{TF} \\ \obl{case\_rest}}
   {\Gamma_{\neg T}\vd \mathit{if}:\tau_{\le}}
  \\ \obl{case\_true}}
 {\Gamma_1\vd \kw{if}\ c\ \kw{then}\ 0\ \kw{else}\ x-y : \tau_{\le}}
\]
The guard premises ($\Gamma_T\vd c:\Prec{\mathsf{true}}$, etc.) are T-CtxSub
followed by T-Var, and are proved. The same pattern derives the angelic
corrected type $x{:}\un{\Nat}{\nu>0}\to y{:}\un{\Nat}{\nu>0}\wand\un{\Nat}{\nu>0}$
for the paper's \textsf{div} guard (\rocq{div\_guard\_program\_tau\_b\_wand}). There, the split is on
whether $y=0$ is reachable.

\subsection{Persistence: two independent generator calls (\S5)}

For the generator form of $e_{\mathit{ND}}$
(\rocq{e\_ND\_gen\_program\_extended\_type}),
\[\kw{let}\ f=\lambda\_.\mathsf{gen}_{\Nat}()\ \kw{in}\ \kw{let}\ a=\mathsf{true}\ \kw{in}\
\kw{let}\ y_1 = f\,a\ \kw{in}\ \kw{let}\ y_2 = f\,a\ \kw{in}\ y_1-y_2 \;:\; \un{\Nat}{\top},\]
the derivation is linear:
\begin{enumerate}
\item T-Let: $f : \pers(\ov{\Bool}{\top}\to\un{\Nat}{\top})$ by \RN{T-PersistIntro}, and
  $a:\Prec{\mathsf{true}}$. Let $\Gamma_a$ be the resulting context.
\item T-Let: $\Gamma_a\vd f\,a:\un{\Nat}{\top}$, using \RN{T-PersE} as the
  subtyping $\pers\tau\sub\tau$ \obl{persist\_elim}.
\item T-CtxSub: $\Gamma_a\entc y_1{:}\un{\Nat}{\top}\ \preceq\ \Gamma_a\sep y_1{:}\un{\Nat}{\top}$
  \qed{comma\_to\_star}. This holds because $\pers$ makes the $f$-projection a
  singleton and $a$ is precise, so the entangled context factors as a product.
\item T-LetD with $\Gamma_1=\Gamma_a$: $\Gamma_a\vd f\,a:\un{\Nat}{\top}$, and
  $y_1{:}\un{\Nat}{\top}\sep y_2{:}\un{\Nat}{\top}\vd y_1-y_2:\un{\Nat}{\top}$ by
  BinOp and Sub \obl{minus\_sub}.
\end{enumerate}
The paper's duplication $x{:}\pers\tau\preceq x{:}\pers\tau\sep x{:}\pers\tau$
is ill-formed in Rocq, where $\sep$ requires disjoint domains. Step~3 replaces
it. With $f=\lambda x.z$ or $f=\lambda x.0\oplus\lambda x.\mathsf{gen}()+1$
instead of the generator, the corresponding judgments are semantically false,
consistent with the paper's \ding{55} marks.

\subsection{Disproving merge by unrolling (Fig.~8)}\label{sec:merge}

Let $\mathsf{O}[l] \triangleq \ov{\List}{\nu=l}$ and $\mathsf{P}[l]\triangleq\Prec{l}$.
We derive (\rocq{merge\_cex\_extended\_type})
\[\vd \mathit{merge} : \mathsf{P}[1;2]\wand \mathsf{P}[3;4]\wand \ov{\List}{\neg\mathit{sorted}(\nu)}\Sum\ov{\List}{\top},\]
the paper's disproof type instantiated at the counterexample inputs, since
$\mathit{merge}\,[1;2]\,[3;4] = [1;3;2;4]$. The derivation has three
unrolling levels; each one uses T-FixUnroll, so T-Fix is never used.
\begin{itemize}
\item \textbf{Level B} (\rocq{mu\_fix00}): $\mathit{fix}:\mathsf{O}[\,]\to\mathsf{O}[\,]\to\mathsf{O}[\,]$.
 The steps are FixUnroll, Lam, Lam, then \RN{ListMatchNilOnly}; the equation
 premise $[\,]:\Eq{xs}$ is Const followed by Sub \obl{nil\_eq}, and the branch
 is T-VarComma.
\item \textbf{Level A} (\rocq{mu\_fix24}): $\mathit{fix}:\mathsf{O}[2]\to\mathsf{O}[4]\to\mathsf{O}[2;4]$.
\[
\inferrule*[right={FixUnroll; Lam$^2$}]
 {\inferrule*[right=CtxSub]
   {\inferrule*[right=ListMatchCons]
     {\mathcal{E}_1 \\ \mathcal{M}_{ys}}
     {G_{A1}\vd \kw{match}\ xs\ \ldots : \mathsf{O}[2;4]}
    \\ \qed{pat_1}}
   {G_{A0} \vd \kw{match}\ xs\ \ldots : \mathsf{O}[2;4]}}
 {\vd \kw{fix}\ \mathit{merge}.\ldots : \mathsf{O}[2]\to\mathsf{O}[4]\to\mathsf{O}[2;4]}
\]
\[
\mathcal{E}_1 =
\inferrule*[right=Sub]
 {\inferrule*[right=ListCons]{\ldots}{G_{A1}\vd h_1::t_1:\Prec{h_1::t_1}} \\ \obl{eq_1}}
 {G_{A1}\vd h_1::t_1:\Eq{xs}}
\]
\[
\mathcal{M}_{ys} =
\inferrule*[right=CtxSub]
 {\inferrule*[right=ListMatchCons]{\mathcal{E}_2\ \obl{eq_2} \\ \mathcal{D}_{\mathit{step}}}
   {G_{A2}\vd \kw{match}\ ys\ \ldots:\mathsf{O}[2;4]}
  \\ \qed{pat_2}}
 {G_{A1}\vd \kw{match}\ ys\ \ldots:\mathsf{O}[2;4]}
\]
 with $G_{A0} = xs{:}\mathsf{O}[2]\entc ys{:}\mathsf{O}[4]$.
 Here $G_{A1} = G_{A0}\entc h_1{:}\ov{\Nat}{\nu{=}2}\entc t_1{:}\mathsf{O}[\,]$
 binds the pattern variables, introduced by the proved pattern lemma
 \qed{pat_1}. $\mathcal{D}_{\mathit{step}}$ continues as follows: a Let binds
 $t = \mathit{fix}\,t_1\,t_2 : \mathsf{O}[\,]$ by AppFun at level B, with the
 partial application typed by T-VarComma; a Let binds
 $\mathit{lt} = h_1{<}h_2:\Prec{h_1<h_2}$; then MatchTrueOnly
 \obl{lt\_true}, a Let for $h_2::t$, and Sub on $h_1::(h_2::t)$
 \obl{leaf}.
\item \textbf{Level O}: the same shape at $\mathsf{P}[1;2]\to(\mathsf{P}[3;4]\wand\ldots)$,
 calling level A for $\mathit{merge}\,[2]\,[4]$. The final Sub to
 $\ov{\List}{\neg\mathit{sorted}}\Sum\ov{\List}{\top}$ is
 \obl{unsorted}: $[1;3;2;4]$ is unsorted. The outer ``$\to$'' becomes
 ``$\wand$'' through the general lemma $\tau_x\to\tau\sub\tau_x\wand\tau$ \qed{arrow\_wand}.
\end{itemize}
Every Admitted obligation here concerns a context in which all variables are
constants, so each fibering has a single fiber.

\subsection{Summary of all case studies}

\begin{center}\small
\begin{tabular}{p{0.24\textwidth}p{0.71\textwidth}}
\toprule
Paper item & Outcome in Rocq\\
\midrule
\S1 $\tau_a$ for \textsf{div}/minus & Derived; the minus variant uses nested CtxCase (\S\ref{sec:minus}).\\
\S1 $\tau_b$ with only ``$\to$'' & \emph{False}: on $\{x{=}n,y{=}m\mid 1\le n\le m\}$, $x-y$ is always $0$. The corrected ``$\wand$'' version is derived. The informal p.2 fix (only $f$'s parameter independent) is insufficient; both p.8 fixes are fully proved.\\
\S2 contexts $C_{\mathit{dis}}$, $C_{\mathit{en}}$, $\Gamma_{\mathit{dis}}$, $\Gamma_{\mathit{en}}$ & All \ding{51}/\ding{55} judgments are confirmed. $x-(z+y)$ is derived with T-LetI by giving $z+y$ a demonic type.\\
\S2/\S3 $\Sum$ examples, Ex.~3.4 & Derived, including the demonic version $x{:}\ov{\Nat}{\top}\vdash\kw{if}\ x{=}0\ldots$ via CtxCase.\\
Ex.~3.1 & The conclusion is derived, but the paper's intermediate premise $x{:}\ov{\Nat}{\nu<10}\entc y{:}\ov{\Nat}{\nu<7}\vdash x-y:\ov{\Nat}{\nu\le3}$ is \emph{invalid} ($(9,0)$). The derivation must keep $y$ precise.\\
Ex.~3.3 (\textsf{sub}) & True, but \emph{not derivable}: the T-Fix issue of \S\ref{sec:open}.\\
p.8 nested $f\,(f\,y)$ & Right-nested A-normal form is not derivable, because \kw{let} entangles $r_2$ with $x$. The monadic-normal-form nesting is derived.\\
\S4 Ex.~4.5, 4.6 & Derived; several fully proved.\\
\S5 generators, $\pers$ & $\mathsf{gen}_{\Bool}():\ov{\Bool}{\nu{=}\mathsf{true}}\Sum\ov{\Bool}{\nu{=}\mathsf{false}}$ derived; the generator form of $e_{\mathit{ND}}$ derived; the non-persistent variants are false, as the paper marks.\\
Fig.~6(a) & (1) false; (2) and (3) derived (\S\ref{sec:fig6a}).\\
Fig.~6(b) & The demonic-$x$ type is false; the angelic version and the safety version are derived.\\
Fig.~7 (\textsf{bst\_gen}) & The stated type is \emph{vacuously true}: $\pers(x{:}\ov{\Nat}{\top}\to y{:}\ov{\Nat}{x<\nu}\to\un{\Nat}{x<\nu<y})$ is uninhabited, because $x{=}1,y{=}2$ gives an empty exact angelic atom. The ``$\le$'' repair is false by totality. A precise-range version is true; its outer layers and range call are derived, and the rest is blocked by T-Fix.\\
Fig.~8 (merge) & The counterexample instance is derived by unrolling (\S\ref{sec:merge}). The literal type with $\un{\List}{\top}$ parameters does not witness the bug: a correct merge also has it, via $\mathit{merge}\,[\,]\,l=l$. The version with sorted angelic parameters is blocked by T-Fix.\\
\bottomrule
\end{tabular}
\end{center}

% =====================================================================
\section{Open issues}\label{sec:open}

\paragraph{1. The recursive binder of T-Fix is vacuous on mixed capabilities.}
\RN{T-Fix} types the body under
$f : (\ov{b}{\phi}\inter\ov{b}{\nu\prec x})\to\tau$. The arrow denotation is a
Kripke implication over whole-world argument extensions, with the argument
type fibered by $x$. On any capability that contains a $\prec$-minimal $x$
(for example $x=0$, or $xs=[\,]$) alongside larger values, no argument satisfies
$\nu\prec x$ in the minimal fiber. The implication therefore holds vacuously,
and \emph{every} $f$ satisfies the binding, including a diverging one. Since
type denotations include totality, the premise of T-Fix is false whenever the
body actually recurses, and no rule can derive it. This was found
independently in Ex.~3.3, Fig.~7 and Fig.~8, and confirmed by audits, but it is
not yet a machine-checked lemma. The paper's arrow denotation seems to share
it. A plausible repair is to fiber the recursive binding by $x$, i.e.\ evaluate
$f$'s type per $x$-fiber ($x\triangleright\ldots$), as a new T-Fix variant
with its own soundness proof. This decision is pending.

\paragraph{2. Unproved semantic obligations.} 129 obligations remain
\rocq{Admitted}. Most are of a few recurring kinds: $\Sum$-splits of contexts,
operator-result subtyping on single or constant fibers, and $\pers$
elimination. They are not proved because the library lacks world-level
introduction and elimination lemmas for $\mathsf{FPlus}$, world filters, and
precise types at $\kw{ret}\,x$. Adding these generically would discharge
most of the obligations at once.

\paragraph{3. Soundness of the extended judgment as a whole.}
\rocq{ExtendedRuleSoundness.v} proves the \emph{new} rules sound in isolation.
The pre-existing convenience rules (\rocq{EXT\_BinOp}, \rocq{EXT\_AppOpValue},
\rocq{EXT\_TreeNode}, \rocq{EXT\_ListCons}, the named
\rocq{Lam}/\rocq{Let}/\rocq{LamD}/\rocq{LetD}/\rocq{Fix} variants) still have no
soundness proof. A full Fundamental theorem for $\vd$, by induction reusing
the per-rule lemmas, is the natural next step.

\paragraph{4. Smaller semantic mismatches with the paper.}
\begin{itemize}
\item The angelic bottom $\un{b}{\bot}$ is uninhabited under exact atoms,
  whereas the paper states $\ov{\Nat}{\nu>0}\sub\un{\Nat}{\bot}$ (p.5).
\item Persistent duplication needs the disjoint-domain workaround of \S4.3.
\item T-AppFun and T-AppOp take only variable arguments, so several programs
  were A-normalized; all such variants are documented in the files.
\end{itemize}

% =====================================================================
\section{Conclusion}

With four rule extensions, each proved sound against the core denotation,
almost every case study of the paper now has a mechanized derivation whose
structural part is checked by Rocq. Making the derivations explicit exposed
several issues in the paper:
\begin{itemize}
\item a dependency bug in Fig.~6(a) setting (1), which the ``$\wand$'' setting
  (2) repairs;
\item a false type in Fig.~6(b);
\item a vacuous type in Fig.~7;
\item a non-witnessing disproof type in Fig.~8;
\item an invalid intermediate judgment in Ex.~3.1;
\item an unfaithful translation of $\un{b}{\top}$.
\end{itemize}
The main remaining obstacle is the vacuity of T-Fix's recursive binder on
mixed capabilities. Resolving it, for example with a fibered recursive
binding, would unblock the three recursive examples: Ex.~3.3, Fig.~7 and the
sorted-input form of Fig.~8.

\end{document}
